% Source presentations for register row N-c-020.

\subsection{Local heap reasoning, framing, and bi-abduction}
\label{sec:p09-presentation-n-c-020}

\subsubsection{Exact spatial satisfaction}

Let a stack \(s\) assign values to program variables and let a heap \(h\) be a
finite partial map from addresses to values.  In addition to ordinary
first-order assertions about \(s\), use the following spatial assertions:
\[
\begin{aligned}
 (s,h)\models\mathsf{emp}
 &\Longleftrightarrow \operatorname{dom}(h)=\varnothing,\\
 (s,h)\models e\mapsto e'
 &\Longleftrightarrow
   \operatorname{dom}(h)=\{\llbracket e\rrbracket_s\}
   \text{ and }
   h(\llbracket e\rrbracket_s)=\llbracket e'\rrbracket_s,\\
 (s,h)\models P*Q
 &\Longleftrightarrow
   \exists h_1,h_2\,
   [h=h_1\uplus h_2\land h_1\perp h_2\land
    (s,h_1)\models P\land(s,h_2)\models Q],\\
 (s,h)\models P\mathbin{-\!*}Q
 &\Longleftrightarrow
   \forall h'\perp h\,
   ((s,h')\models P\to(s,h\uplus h')\models Q).
\end{aligned}
\]
Here \(\mathsf{emp}\), points-to, separating conjunction, and separating
implication are assertion constructors.  They are needed because ordinary
conjunction evaluates both conjuncts on the same entire heap and cannot state
that two owned cells are disjoint.  Their standing is semantic rather than
operational, they differ by their exact heap-splitting or heap-extension
clauses, and the first three are used in the local update and frame derivation
below \cite[\S\S2--3]{OHearnReynoldsYang2001}
\cite[pp.~55--61]{ReynoldsSeparation2002}.

Two eliminable abbreviations make the four primitive axioms legible:
\[
 e\mapsto-\ \equiv\ \exists w\,(e\mapsto w),
 \qquad
 e\mapsto(v_1)\ \equiv\ e\mapsto v_1,
 \qquad
 e\mapsto(v_1,\ldots,v_m)
 \ \equiv\
 (e\mapsto v_1)*((e+1)\mapsto v_2)*
 \cdots*((e+m-1)\mapsto v_m)\quad(m\geq2).
\]
The first asserts ownership of one cell with unspecified contents.  The
second asserts ownership of \(m\) consecutive cells with the displayed
contents; it is not a new heap constructor, and differs from singleton
points-to only by expanding to the displayed finite separating conjunction.
Both abbreviations are needed and used in the primitive rules and list proof
below.

For \(m\geq1\), the following local triples are valid:
\[
\begin{array}{rcll}
 \{e\mapsto u\}
   &x:=[e]&
 \{x=u\land e\mapsto u\},
 &x\notin\operatorname{FV}(e),\\[2pt]
 \{e\mapsto-\}
   &[e]:=v&
 \{e\mapsto v\},&\\[2pt]
 \{\mathsf{emp}\}
   &x:=\mathsf{cons}(v_1,\ldots,v_m)&
 \{x\mapsto(v_1,\ldots,v_m)\},
 &x\notin\bigcup_{j=1}^{m}\operatorname{FV}(v_j),\\[2pt]
 \{e\mapsto-\}
   &\mathsf{dispose}\ e&
 \{\mathsf{emp}\}.&
\end{array}
\]
The parameter \(u\) is a rigid logical value: it is supplied when the
assertion is evaluated, is not a stack or heap component, is unchanged by the
command, and the triple is uniform in its value.  This standing is needed so
lookup can relate the new stack value of \(x\) to the unchanged heap cell
without treating \(u\) as an assignable program variable.  The side condition
on lookup keeps the address expression unchanged when \(x\) is assigned; the allocation
condition keeps the initializing expressions independent of the newly chosen
address.  These are fault-avoiding partial-correctness triples.  In the first,
second, and fourth rows the singleton precondition makes the accessed address
active, so the operational command cannot fault and its one transition gives
the displayed postcondition.  In the allocation row, the operational rule
chooses a fresh consecutive block, initializes it with the old values of the
\(v_j\), and stores its first address in \(x\).  The frame rule, rather than a
larger primitive axiom, reattaches any disjoint heap
\cite[pp.~61--63]{ReynoldsSeparation2002}.  The displayed lookup row is the
source-valid specialization in which the assigned variable is absent from the
address expression.  Reynolds also prints a more general auxiliary-variable
schema which retains the old value of \(x\); that schema is not silently
identified with this restricted row.

\subsubsection{Frame soundness and one symbolic update}

The registered frame rule is
\[
 \frac{\{P\}\ C\ \{Q\}}
      {\{P*R\}\ C\ \{Q*R\}}.
\]
Its side condition is that no stack variable modified by \(C\) occurs free in
\(R\).  Its proof uses command locality from
Definition~\ref{def:p09-separation-logic}.  Split the initial heap as
\(h=h_P\uplus h_R\).  Safety monotonicity permits execution in the larger heap;
the frame property explains every terminating result as
\(h'_P\uplus h_R\), with the framed heap unchanged.  The premise gives \(Q\) on
\(h'_P\), while the free-variable side condition preserves the stack facts used
by \(R\).  Hence the result satisfies \(Q*R\).  Disjointness without locality or
the variable side condition is insufficient \cite[\S\S3--5]{OHearnReynoldsYang2001}.

The stack-variable side condition has its own one-line counterexample.  The
local heap triple
\[
 \{\mathsf{emp}\}\ x:=0\ \{\mathsf{emp}\}
\]
is valid.  If the frame rule admitted
\(R\equiv x=1\land\mathsf{emp}\), it would infer
\[
 \{\mathsf{emp}*R\}\ x:=0\ \{\mathsf{emp}*R\}.
\]
An empty-heap state with \(x=1\) satisfies the precondition by splitting the
heap into two empty parts, but the final stack has \(x=0\) and falsifies the
postcondition.  Heap locality therefore does not make a framed assertion
independent of modified stack variables.

For a concrete symbolic step, suppose \(x,y,z\) denote values with distinct
addresses for \(x\) and \(y\).  The small axiom and frame rule derive
\[
 \{x\mapsto y * y\mapsto0\}\ [x]:=z\quad
 \{x\mapsto z * y\mapsto0\}.
\]
The update consumes and re-establishes the \(x\)-cell; \(y\mapsto0\) is carried
unchanged.  If \(x=y\), the precondition itself is false because one singleton
heap cannot be split into two disjoint cells at the same address.  This is the
shortest explicit aliasing boundary: an ordinary conjunction would not supply
the separation required by the proof.

Reynolds's 2002 article gives the stack--heap interpretation, points-to,
separating conjunction and implication, the frame rule, and list-manipulation
derivations.  It is a semantic and proof-calculus source; it does not prove that
an automatic analyzer discovers every valid specification
\cite{ReynoldsSeparation2002}.

\subsubsection{A complete local list-reversal step}

Let \(\mathsf{nil}\) be a distinguished value which is not an address, let
\(\epsilon\) be the empty finite sequence, let \(a\cdot\alpha\) prefix \(a\)
to \(\alpha\), and let \(\alpha^\dagger\) reverse \(\alpha\).  Define the
assertion family
\[
\begin{aligned}
 \operatorname{list}_{\epsilon}(i,j)
   &\equiv \mathsf{emp}\land i=j,\\
 \operatorname{list}_{a\cdot\alpha}(i,k)
   &\equiv
   \exists j\,
   \bigl(i\mapsto(a,j)*
         \operatorname{list}_{\alpha}(j,k)\bigr).
\end{aligned}
\]
This family is needed to describe an unbounded linked segment while retaining
its abstract sequence of data.  It stands as a recursively defined family of
heap assertions, not as a command or a second heap model.  It differs from one
points-to assertion by owning every record in the segment, and from bare
reachability by recording the sequence and enforcing disjoint ownership.  It
is used in the following proof.

Starting from \(\operatorname{list}_{\alpha_0}(i,\mathsf{nil})\), consider
\[
\begin{split}
 j&:=\mathsf{nil};\\
 \mathbf{while}\ i\neq\mathsf{nil}\ \mathbf{do}\quad
   k:=[i+1];\ [i+1]:=j;\ j:=i;\ i:=k .
\end{split}
\]
The loop assertion is
\[
 \mathcal I\equiv
 \exists\alpha,\beta\,
 \left[
  \left(
   \operatorname{list}_{\alpha}(i,\mathsf{nil})*
   \operatorname{list}_{\beta}(j,\mathsf{nil})
  \right)
  \land
  \alpha_0^\dagger=\alpha^\dagger\cdot\beta
 \right].
\]
It is established after \(j:=\mathsf{nil}\) by taking
\(\alpha=\alpha_0\) and \(\beta=\epsilon\).

For preservation, \(\mathcal I\land i\neq\mathsf{nil}\) permits the first list
to be unfolded with \(\alpha=a\cdot\gamma\).  After the lookup, the owned heap
and pure sequence equation have the form
\[
 \exists a,\gamma,\beta\,
 \left[
  i\mapsto(a,k)*
  \operatorname{list}_{\gamma}(k,\mathsf{nil})*
  \operatorname{list}_{\beta}(j,\mathsf{nil})
  \ \land\
  \alpha_0^\dagger=\gamma^\dagger\cdot a\cdot\beta
 \right].
\]
The mutation \([i+1]:=j\) changes only the second cell of the \(i\)-record:
\[
 \exists a,\gamma,\beta\,
 \left[
  i\mapsto(a,j)*
  \operatorname{list}_{\gamma}(k,\mathsf{nil})*
  \operatorname{list}_{\beta}(j,\mathsf{nil})
  \ \land\
  \alpha_0^\dagger=\gamma^\dagger\cdot a\cdot\beta
 \right].
\]
Folding the first and third factors gives
\[
 \exists a,\gamma,\beta\,
 \left[
  \operatorname{list}_{\gamma}(k,\mathsf{nil})*
  \operatorname{list}_{a\cdot\beta}(i,\mathsf{nil})
  \ \land\
  \alpha_0^\dagger=\gamma^\dagger\cdot(a\cdot\beta)
 \right].
\]
The assignments \(j:=i;i:=k\) therefore re-establish \(\mathcal I\), now with
\(\gamma\) as the unreversed segment and \(a\cdot\beta\) as the reversed
segment.  On exit \(i=\mathsf{nil}\).  Since \(\mathsf{nil}\) is not an
address, the first list assertion can then use only its empty-sequence clause;
hence \(\alpha=\epsilon\), \(\beta=\alpha_0^\dagger\), and
\[
 \operatorname{list}_{\alpha_0^\dagger}(j,\mathsf{nil})
\]
holds.  In full, the derivation proves
\[
 \{\operatorname{list}_{\alpha_0}(i,\mathsf{nil})\}\,
 j:=\mathsf{nil};
 \mathbf{while}\ i\neq\mathsf{nil}\ \mathbf{do}\quad
   (k:=[i+1];[i+1]:=j;j:=i;i:=k)\,
 \{\operatorname{list}_{\alpha_0^\dagger}(j,\mathsf{nil})\}.
\]
Each heap step used one small axiom plus framing of the other two
disjoint factors \cite[pp.~63--65]{ReynoldsSeparation2002}.  A further list
rooted at an unmodified stack variable can be framed around the whole proof;
no pairwise reachability side conditions need be re-proved.  The displayed
triple is fault-avoiding partial correctness.  A total-correctness proof would
add the length of the unreversed sequence \(\alpha\) as a separate decreasing
variant.

\subsubsection{Bi-abduction and its conditional guarantee}

Calcagno et al. ask for assertions \(M\) and \(F\) satisfying
\[
 A*M\ \vdash\ G*F.
\]
The unknown \(M\) is the missing precondition or anti-frame; \(F\) is the
residual frame.  These names are necessary because one assertion is added to
make a desired symbolic step possible while the other records resources left
over.  Both stand as assertions solved for in one entailment problem, differ by
which side of the entailment they occupy, and are used by compositional summary
inference \cite[\S\S2--3]{CalcagnoDistefanoOHearnYang2011}.

For instance, let
\[
 A=(x\mapsto y)*(z\mapsto0),
 \qquad
 G=(x\mapsto y)*(y\mapsto0),
\]
under the explicit assumption that the three displayed addresses are pairwise
distinct.  One solution is
\[
 M=y\mapsto0,
 \qquad
 F=z\mapsto0,
\]
because associativity and commutativity of disjoint heap union make
\(A*M\) and \(G*F\) denote the same three-cell heaps.  The example explains the
two inferred roles but does not assert uniqueness or completeness of the
solver.

The same sequent becomes a path summary rather than remaining a static heap
puzzle.  Put \(C\equiv[y]:=1\) and
\(G'\equiv(x\mapsto y)*(y\mapsto1)\).  The small mutation axiom, framed by the
\(x\)-cell, proves \(\{G\}\,C\,\{G'\}\).  Framing once more gives
\(\{G*F\}\,C\,\{G'*F\}\), while the solved entailment supplies
\(A*M\vdash G*F\).  Consequence therefore yields the inferred one-path
summary
\[
 \{A*M\}\ C\ \{G'*F\}.
\]
Thus \(M=y\mapsto0\) is accumulated into the required entry heap and
\(F=z\mapsto0\) is propagated unchanged to the exit heap.  This derivation
states what a sound bi-abduction answer licenses; it does not prescribe a
search algorithm for finding \(M,F\)
\cite[\S\S3--4]{CalcagnoDistefanoOHearnYang2011}.

The program-analysis soundness result is conditional on termination of the
inference procedure and soundness of the primitive rules: an inferred summary
denotes a true fault-avoiding triple.  The source does not claim that every true
summary is inferred, and its final analysis theorem delegates operational-
semantics details to the cited framework
\cite[thm.~4.14]{CalcagnoDistefanoOHearnYang2011}.

\paragraph{Historical and tools boundary.}
O'Hearn's 2019 retrospective traces the path from small footprints and framing
to bi-abductive compositional analysis and deployed tools.  It supports that
historical and engineering account, not a replacement formal semantics or a
new completeness theorem \cite{OHearn2019}.  Successful deployment is evidence
of scalability on the reported programs; it is not evidence that every heap
error is found.
